% =============================================================================
% Section: ATLAS Recreation (MCQ)
% Bib: refs/mcq_atlas.bib
%
% FIGURES REFERENCED (\includegraphics{figures/<basename>}).
% Integrator: copy each source file below into the paper's figures/ directory.
%   figures/atlas_arc_heldout_se0.2.png
%     <- /Users/arhant/Documents/EDLM/olmo-eval-full/AdaptiveTesting/Research/01_MCQ_ATLAS/figures/atlas_arc_heldout_se0.2.png
%   figures/atlas_arc_heldout_se0.3.png
%     <- /Users/arhant/Documents/EDLM/olmo-eval-full/AdaptiveTesting/Research/01_MCQ_ATLAS/figures/atlas_arc_heldout_se0.3.png
%   figures/atlas_arc_0p5_7b_heldout_se0.2.png
%     <- /Users/arhant/Documents/EDLM/olmo-eval-full/AdaptiveTesting/Research/01_MCQ_ATLAS/figures/atlas_arc_0p5_7b_heldout_se0.2.png
%   figures/atlas_arc_0p5_7b_heldout_se0.3.png
%     <- /Users/arhant/Documents/EDLM/olmo-eval-full/AdaptiveTesting/Research/01_MCQ_ATLAS/figures/atlas_arc_0p5_7b_heldout_se0.3.png
%   figures/atlas_selfresp_arc_se0.2.png
%     <- /Users/arhant/Documents/EDLM/olmo-eval-full/AdaptiveTesting/Research/01_MCQ_ATLAS/data/atlas_selfresp/figures/atlas_selfresp_arc_se0.2.png
%   figures/atlas_selfresp_arc_se0.3.png
%     <- /Users/arhant/Documents/EDLM/olmo-eval-full/AdaptiveTesting/Research/01_MCQ_ATLAS/data/atlas_selfresp/figures/atlas_selfresp_arc_se0.3.png
%   figures/atlasrep_mae_vs_se.png
%     <- /Users/arhant/Documents/EDLM/olmo-eval-full/AdaptiveTesting/Research/01_MCQ_ATLAS/figures/atlasrep_mae_vs_se.png
%   figures/atlasrep_items_vs_se.png
%     <- /Users/arhant/Documents/EDLM/olmo-eval-full/AdaptiveTesting/Research/01_MCQ_ATLAS/figures/atlasrep_items_vs_se.png
%   figures/atlasrep_theta_mae_vs_se.png
%     <- /Users/arhant/Documents/EDLM/olmo-eval-full/AdaptiveTesting/Research/01_MCQ_ATLAS/figures/atlasrep_theta_mae_vs_se.png
%   figures/atlasrep_ifeval_pirt_vs_actual_se_0.2.png
%     <- /Users/arhant/Documents/EDLM/olmo-eval-full/AdaptiveTesting/Research/01_MCQ_ATLAS/figures/atlasrep_ifeval_pirt_vs_actual_se_0.2.png
%   figures/atlasrep_gpqa_pirt_vs_actual_se_0.2.png
%     <- /Users/arhant/Documents/EDLM/olmo-eval-full/AdaptiveTesting/Research/01_MCQ_ATLAS/figures/atlasrep_gpqa_pirt_vs_actual_se_0.2.png
%   figures/atlasrep_math_pirt_vs_actual_se_0.2.png
%     <- /Users/arhant/Documents/EDLM/olmo-eval-full/AdaptiveTesting/Research/01_MCQ_ATLAS/figures/atlasrep_math_pirt_vs_actual_se_0.2.png
%   figures/atlasrep_bbh_pirt_vs_actual_se_0.2.png
%     <- /Users/arhant/Documents/EDLM/olmo-eval-full/AdaptiveTesting/Research/01_MCQ_ATLAS/figures/atlasrep_bbh_pirt_vs_actual_se_0.2.png
%   figures/atlasrep_musr_pirt_vs_actual_se_0.2.png
%     <- /Users/arhant/Documents/EDLM/olmo-eval-full/AdaptiveTesting/Research/01_MCQ_ATLAS/figures/atlasrep_musr_pirt_vs_actual_se_0.2.png
%
% PREAMBLE REQUIREMENTS (main document): \usepackage{graphicx}, \usepackage{xcolor},
%   \usepackage{booktabs}. If \url is used in the .bib, load \usepackage{hyperref} or {url}.
%
% Convention: human prose (from the outline) is left in the default color.
%   AI-added prose is wrapped in \aiadd{...} (renders blue). Fully-AI subsections
%   close with an \aibullets{...} takeaway list.
% =============================================================================

\providecommand{\aiadd}[1]{\textcolor{blue}{#1}}
\providecommand{\aibullets}[1]{{\color{blue}#1}}

\section{ATLAS Recreation}
\label{sec:mcq-atlas-recreation}

% --- HUMAN prose (outline PDF, lightly cleaned) -------------------------------
While ATLAS showcases promising results for reduction in benchmark size, there is
a gap between the ATLAS results and actual applications for evaluations during
model training. The ability to apply IRT calibration to a benchmark without public
response data requires a proper analysis of the upfront cost required to reach a
desired correlation. The ability for the benchmarks to discriminate between models
at parameter ranges without many calibration models was also not shown explicitly.
% --- end HUMAN prose ----------------------------------------------------------

\aiadd{To probe that gap we rebuild the ATLAS pipeline~\cite{li2026adaptive} on
our own data and ask four questions. Does a published item bank transfer to a
different pool of models? Does refitting the bank on a narrow parameter range help
or hurt? How much of any transfer loss is a scoring-protocol mismatch rather than
a limit of item response theory? And do the error patterns ATLAS reported on its
five benchmarks reappear when the same analysis runs on the newer Open LLM
Leaderboard tasks~\cite{fourrier2024openllm}? Throughout, every correlation is the
Pearson $r$ between the CAT-predicted benchmark score and the true full-benchmark
accuracy on held-out models.}

% =============================================================================
\subsection{A published bank transfers; a range-restricted refit does not}
\label{subsec:atlas-transfer}

\aiadd{We run ATLAS-style adaptive testing on ARC~\cite{clark2018arc} two ways and
score both on the same 60 held-out models. The first condition uses the published
ATLAS 3PL bank, which was calibrated on thousands of models across a wide parameter
range. The second refits the identical pipeline on only the 0.5 to 7B slice of
models, about 1{,}686 of them, with roughly 95\% at 7B. Table~\ref{tab:atlas-transfer}
reports both at two stopping thresholds.}

\aiadd{The published bank recovers held-out ARC accuracy at $r=0.83$ using about 21
of its 1{,}172 items at $\mathrm{SE}\le0.2$ (Figure~\ref{fig:atlas-transfer}, top
row). Refitting on the narrow 0.5 to 7B slice drops the correlation to $r=0.59$ at
the same threshold, even though the mean absolute error is almost unchanged
(Figure~\ref{fig:atlas-transfer}, bottom row). A range-matched bank is cheaper to
build, but it discriminates worse when the calibration pool is small and
concentrated at a single model size. The loss shows up in ranking correlation, not
in average error.}

\begin{table}[t]
  \centering
  \caption{\aiadd{ATLAS-style CAT on ARC, scored on the same 60 held-out models.
  A published wide-pool 3PL bank versus the same pipeline refit on only the 0.5 to
  7B model slice.}}
  \label{tab:atlas-transfer}
  \small
  \begin{tabular}{l c c c c}
    \toprule
    Calibration pool & SE stop & Pearson $r$ & MAE & Mean items \\
    \midrule
    Published ATLAS bank  & 0.2 & 0.830 & 0.147 & 21.1 \\
    Published ATLAS bank  & 0.3 & 0.738 & 0.172 & 9.9  \\
    Refit, 0.5 to 7B      & 0.2 & 0.589 & 0.147 & 12.5 \\
    Refit, 0.5 to 7B      & 0.3 & 0.585 & 0.150 & 8.3  \\
    \bottomrule
  \end{tabular}
\end{table}

\begin{figure*}[t]
  \centering
  \includegraphics[width=0.48\linewidth]{figures/atlas_arc_heldout_se0.2.png}\hfill
  \includegraphics[width=0.48\linewidth]{figures/atlas_arc_heldout_se0.3.png}\\[3pt]
  \includegraphics[width=0.48\linewidth]{figures/atlas_arc_0p5_7b_heldout_se0.2.png}\hfill
  \includegraphics[width=0.48\linewidth]{figures/atlas_arc_0p5_7b_heldout_se0.3.png}
  \caption{\aiadd{p-IRT predicted accuracy against actual full-benchmark accuracy on
  60 held-out ARC models. Top row: the published ATLAS 3PL bank at $\mathrm{SE}\le0.2$
  (left) and $\mathrm{SE}\le0.3$ (right). Bottom row: the same pipeline refit on only
  the 0.5 to 7B slice. The wide-pool bank tracks the $y=x$ line more tightly.}}
  \label{fig:atlas-transfer}
\end{figure*}

\aibullets{%
\begin{itemize}
  \item The published wide-pool bank transfers to a fresh set of models at $r=0.83$
  ($\mathrm{SE}\le0.2$) using about 21 of 1{,}172 items.
  \item A bank refit only on 0.5 to 7B models falls to $r=0.59$, so parameter-range
  coverage in the calibration pool matters more than matching the target range.
  \item MAE stays near 0.147 in both conditions, so the loss is in ranking
  correlation rather than average error.
\end{itemize}}

% =============================================================================
\subsection{Prompt-matched responses recover near-ceiling transfer}
\label{subsec:atlas-selfresp}

\aiadd{The transfer test above carries a confound. Our 0-shot loglikelihood
responses were scored against a bank that ATLAS calibrated on 25-shot leaderboard
responses, so the prompting protocol does not match. To separate that from any real
IRT transfer error, we replay the same 3PL CAT on ATLAS's own ARC test response
matrix, aligned to the bank by item index. The held-out set here is a seed-7 sample
of 60 models drawn from ATLAS's 417-model ARC test set, which spans a wider
parameter range than our 0.2 to 7B pool, and 650 bank items align to that matrix.}

\aiadd{With prompt-matched responses the published bank reaches $r=0.91$ in under 10
items, and the mean absolute error drops by about 5x, from 0.147 to 0.031 at
$\mathrm{SE}\le0.2$ (Table~\ref{tab:atlas-selfresp}, Figure~\ref{fig:atlas-selfresp}).
Most of the degradation seen in Section~\ref{subsec:atlas-transfer} came from the
scoring-protocol mismatch, not from a failure of the item bank to carry over.}

\begin{table}[t]
  \centering
  \caption{\aiadd{The same published ATLAS bank and 3PL CAT, scored on our 0-shot
  responses versus ATLAS's own (prompt-matched) ARC test responses.}}
  \label{tab:atlas-selfresp}
  \small
  \begin{tabular}{l c c c c}
    \toprule
    Held-out responses   & SE stop & Pearson $r$ & MAE & Mean items \\
    \midrule
    Our 0-shot responses & 0.2 & 0.830 & 0.147 & 21.1 \\
    Our 0-shot responses & 0.3 & 0.738 & 0.172 & 9.9  \\
    ATLAS own responses  & 0.2 & 0.915 & 0.031 & 13.0 \\
    ATLAS own responses  & 0.3 & 0.906 & 0.034 & 9.2  \\
    \bottomrule
  \end{tabular}
\end{table}

\begin{figure}[t]
  \centering
  \includegraphics[width=0.49\linewidth]{figures/atlas_selfresp_arc_se0.2.png}\hfill
  \includegraphics[width=0.49\linewidth]{figures/atlas_selfresp_arc_se0.3.png}
  \caption{\aiadd{The published ATLAS bank replayed on ATLAS's own ARC test
  responses at $\mathrm{SE}\le0.2$ (left) and $\mathrm{SE}\le0.3$ (right). Removing
  the 0-shot versus 25-shot scoring mismatch raises the correlation to $r=0.91$ and
  cuts MAE about 5x relative to Table~\ref{tab:atlas-transfer}.}}
  \label{fig:atlas-selfresp}
\end{figure}

\aibullets{%
\begin{itemize}
  \item On prompt-matched responses the bank recovers accuracy at $r=0.91$
  ($\mathrm{SE}\le0.2$) in about 13 items, and $r=0.906$ at $\mathrm{SE}\le0.3$ in
  about 9.
  \item MAE falls about 5x against the 0-shot replay (0.031 versus 0.147), so a
  matched scoring protocol, not IRT itself, drove most of the earlier error.
\end{itemize}}

% =============================================================================
\subsection{Replicating the ATLAS error analyses on Open LLM Leaderboard tasks}
\label{subsec:atlas-replication}

\aiadd{ATLAS reports a fixed set of error views: a scatter of p-IRT predicted
accuracy against actual accuracy, a histogram of the prediction error, absolute
error against the number of items administered, and error against actual accuracy.
We reproduce this set on five Open LLM Leaderboard v2 tasks~\cite{fourrier2024openllm}
that ATLAS never covered: IFEval~\cite{zhou2023ifeval}, MATH~\cite{hendrycks2021math},
GPQA~\cite{rein2024gpqa}, BBH~\cite{suzgun2023bbh}, and MuSR~\cite{sprague2024musr}.
We use the same frozen 3PL banks and the Fisher-information p-IRT selector, with a
single calibration on a fixed 90/10 split at seed 7, at ATLAS's own SE thresholds of
0.1, 0.2, and 0.3.}

\aiadd{Table~\ref{tab:atlas-replication} gives the per-benchmark correlation,
accuracy MAE, ability (theta) MAE, and the adaptive-versus-random item counts. On
IFEval, GPQA, and MATH the accuracy MAE (0.016 to 0.070) sits in the same 0.02 to
0.05 band ATLAS reports, and it falls as the SE target tightens
(Figure~\ref{fig:atlasrep-curves}, left). Fisher-information selection reaches a
given SE in 1.6x to 5.6x fewer items than a random-item baseline, which is the core
ATLAS efficiency result (Figure~\ref{fig:atlasrep-curves}, center). Ability recovery
on IFEval and MATH (theta MAE 0.07 to 0.21) brackets ATLAS's 0.06 to 0.18 range
(Figure~\ref{fig:atlasrep-curves}, right).}

\aiadd{One preprocessing step decides these numbers. We drop items with
non-positive discrimination ($a\le0$), which are mirt calibration failures or
reverse-scored items, matching the canonical GPQA loader. The GPQA, MuSR, and BBH
banks contain 51\%, 43\%, and 31\% such items. Keeping them collapses GPQA to
$r\approx0$; dropping them recovers $r=0.74$ at $\mathrm{SE}\le0.3$. Under the same
filter MuSR recovers to about $r=0.75$, so it is not the clean negative control it
first appeared to be. BBH stays the weakest case (Figure~\ref{fig:atlasrep-scatter}):
its p-IRT estimates are rank-correlated with actual accuracy ($r\approx0.68$) but
biased, because a single-factor 3PL underfits its multi-subtask mixture.}

\begin{table*}[t]
  \centering
  \caption{\aiadd{ATLAS error analyses reproduced on Open LLM Leaderboard v2 tasks
  (about 110 held-out models each; math 90). Ability MAE is measured against the
  full-bank EAP ability. The last column gives mean items to reach
  $\mathrm{SE}\le0.3$ under adaptive versus random item selection.}}
  \label{tab:atlas-replication}
  \small
  \begin{tabular}{l ccc ccc ccc c}
    \toprule
    & \multicolumn{3}{c}{Pearson $r$} & \multicolumn{3}{c}{Accuracy MAE}
    & \multicolumn{3}{c}{Ability (theta) MAE} & Items @ $\mathrm{SE}\le0.3$ \\
    \cmidrule(lr){2-4}\cmidrule(lr){5-7}\cmidrule(lr){8-10}
    Benchmark & 0.1 & 0.2 & 0.3 & 0.1 & 0.2 & 0.3 & 0.1 & 0.2 & 0.3 & adaptive / random \\
    \midrule
    IFEval & 0.955 & 0.927 & 0.921 & 0.034 & 0.047 & 0.050 & 0.070 & 0.160 & 0.207 & 18 / 62  \\
    GPQA   & 0.710 & 0.757 & 0.737 & 0.034 & 0.066 & 0.070 & 0.104 & 0.357 & 0.407 & 13 / 46  \\
    MATH   & 0.948 & 0.892 & 0.878 & 0.016 & 0.023 & 0.024 & 0.097 & 0.168 & 0.208 & 77 / 148 \\
    MuSR   & 0.893 & 0.770 & 0.746 & 0.075 & 0.091 & 0.095 & 0.362 & 0.629 & 0.707 & 14 / 76  \\
    BBH    & 0.695 & 0.677 & 0.674 & 0.109 & 0.112 & 0.114 & 0.711 & 0.719 & 0.684 & 8 / 13   \\
    \bottomrule
  \end{tabular}
\end{table*}

\begin{figure*}[t]
  \centering
  \includegraphics[width=0.32\linewidth]{figures/atlasrep_mae_vs_se.png}\hfill
  \includegraphics[width=0.32\linewidth]{figures/atlasrep_items_vs_se.png}\hfill
  \includegraphics[width=0.32\linewidth]{figures/atlasrep_theta_mae_vs_se.png}
  \caption{\aiadd{Cross-benchmark error curves reproduced from ATLAS on the Open LLM
  Leaderboard v2 tasks. Left: accuracy MAE against the SE stopping threshold. Center:
  mean items administered, adaptive against a random-item baseline. Right: ability
  (theta) MAE against SE. Tighter SE lowers MAE and raises the item count.}}
  \label{fig:atlasrep-curves}
\end{figure*}

\begin{figure*}[t]
  \centering
  \includegraphics[width=0.32\linewidth]{figures/atlasrep_ifeval_pirt_vs_actual_se_0.2.png}\hfill
  \includegraphics[width=0.32\linewidth]{figures/atlasrep_gpqa_pirt_vs_actual_se_0.2.png}\hfill
  \includegraphics[width=0.32\linewidth]{figures/atlasrep_math_pirt_vs_actual_se_0.2.png}\\[3pt]
  \includegraphics[width=0.32\linewidth]{figures/atlasrep_bbh_pirt_vs_actual_se_0.2.png}\hfill
  \includegraphics[width=0.32\linewidth]{figures/atlasrep_musr_pirt_vs_actual_se_0.2.png}\hfill
  \begin{minipage}[b]{0.32\linewidth}~\end{minipage}
  \caption{\aiadd{Per-benchmark p-IRT predicted accuracy against actual accuracy at
  $\mathrm{SE}\le0.2$ for IFEval, GPQA, MATH (top) and BBH, MuSR (bottom). IFEval and
  MATH track the $y=x$ line closely; BBH is rank-correlated but biased on its large
  multi-subtask bank.}}
  \label{fig:atlasrep-scatter}
\end{figure*}

\aibullets{%
\begin{itemize}
  \item On IFEval, MATH, and GPQA the p-IRT accuracy MAE (0.016 to 0.070) matches
  ATLAS's 0.02 to 0.05 band and improves as the SE target tightens.
  \item Adaptive selection reaches the same SE in 1.6x to 5.6x fewer items than
  random, reproducing the ATLAS efficiency claim on benchmarks ATLAS never covered.
  \item Dropping non-positive-discrimination items is essential: GPQA moves from
  $r\approx0$ to $r=0.74$, and MuSR recovers to about $r=0.75$, so its near-zero
  result elsewhere was a bad-item artifact.
  \item BBH is the weakest replication ($r\approx0.68$, MAE $\approx0.11$): a
  unidimensional 3PL underfits a multi-subtask benchmark.
\end{itemize}}
